% =========================================================
% SECCIÓN 2.3 - SENSORES Y ARQUITECTURA DE ADQUISICIÓN
% =========================================================

La capa de percepción del sistema híbrido está compuesta por un nodo IoT de bajo costo que integra sensores de infrarrojo no dispersivo (NDIR) para \ce{CO2} y \ce{CH4}, sensores ambientales auxiliares (temperatura y humedad relativa) y un microcontrolador con capacidad de procesamiento en el borde. Esta arquitectura responde al principio de ``soberanía de datos'' y diseño frugal: el procesamiento crítico ocurre localmente, minimizando la dependencia de conectividad continua y preservando la trazabilidad del dato crudo \parencite{vanIntegrationInternetofThingsSustainable2022, duobieneDevelopmentWirelessSensor2022}. A continuación se describen los componentes hardware, sus especificaciones metrológicas y la lógica de telemetría.

% ---------------------------------------------------------
\subsection{Sensor NDIR para \texorpdfstring{\ce{CO2}}{CO2} y \texorpdfstring{\ce{CH4}}{CH4}}
\label{subsec:cap2_sensor_ndir}

El dióxido de carbono se midió mediante un sensor NDIR de bajo costo (modelo MH-Z19B, Winsen, China), basado en el principio de absorción selectiva de radiación infrarroja a \SI{4.26}{\micro\metre}, longitud de onda correspondiente a la banda de vibración asimétrica de la molécula de \ce{CO2}. Este dispositivo ofrece un rango nominal de \num{0}--\num{5000}~ppm, resolución de \num{1}~ppm y precisión de $\pm$\,\SI{50}{ppm} $\pm$\,\SI{5}{\percent} de lectura en condiciones de referencia (\SI{25}{\celsius}, HR \SI{50}{\percent}), alimentado a \SI{5}{\volt} DC \parencite{biagiDevelopmentMachineLearningbased2024}. Para el metano, dado que los sensores NDIR de bajo costo con especificidad para \ce{CH4} son menos accesibles comercialmente, se empleó un sensor de semiconductor metal-óxido (MQ-4) calibrado por correlación cruzada contra patrones de metano certificados, operando en rango \num{300}--\num{10000}~ppm con resolución de \num{10}~ppm.

La elección de sensores comerciales de bajo costo implica una deriva instrumental significativa ante fluctuaciones de temperatura y humedad. La absorción infrarroja en la cavidad NDIR es sensible a la densidad de moléculas de \ce{CO2}, que depende tanto de la concentración molar como de la temperatura del gas muestra. A \SI{30}{\celsius} y HR $>$ \SI{70}{\percent} ---condiciones típicas del reactor de bioconversión---, la condensación de vapor de agua sobre la ventana óptica puede atenuar la señal y generar lecturas erráticas \parencite{herrinCollateralDataQuality2021}. Para mitigar este efecto, se implementó un preacondicionamiento físico consistente en una línea de muestreo de teflón (\SI{4}{\milli\metre} diámetro, \SI{30}{\centi\metre} longitud) con membrana hidrofóbica (PTFE, \SI{0.45}{\micro\metre}) y un pequeño desecante de silica gel intercambiable, sin contacto directo con el flujo principal para evitar pérdida de \ce{CO2} por disolución.

% ---------------------------------------------------------
\subsection{Variables ambientales: temperatura y humedad relativa}
\label{subsec:cap2_sensor_ambiental}

La temperatura del sustrato ($T_s$) y la humedad relativa del aire de cabeza libre (HR) se registraron mediante sensores digitales SHT30 (Sensirion, Suiza), con precisión de $\pm$\,\SI{0.3}{\celsius} y $\pm$\,\SI{2}{\percent} HR, respectivamente, comunicados por bus I2C al microcontrolador central. Se desplegaron tres sondas por reactor: una inmersa en el sustrato a \SI{5}{\centi\metre} de profundidad (zona de máxima actividad larval), otra en el aire de cabeza libre y una tercera en contacto con la pared interna del reactor para detectar gradientes térmicos por el metabolismo exotérmico del conjunto larva-sustrato \parencite{bekkerImpactSubstrateMoisture2021}.

La temperatura del sustrato constituye una variable de control crítica porque modula directamente la tasa de asimilación y la respiración larval, siguiendo una cinética tipo Arrhenius o Sharpe--Schoolfield dentro del rango \SI{25}{\celsius}--\SI{35}{\celsius} \parencite{eriksenDynamicModellingFeed2022}. La humedad relativa, por su parte, condiciona la formación de microambientes anaeróbicos: valores $>$ \SI{80}{\percent} favorecen la saturación de poros y la generación de \ce{CH4}, mientras que valores $<$ \SI{50}{\percent} reducen la tasa de asimilación y aumentan la mortalidad \parencite{bekkerImpactSubstrateMoisture2021, salamEffectDifferentEnvironmental2022}. Por ello, las lecturas de HR se utilizaron no solo como variables de compensación cruzada para el sensor NDIR, sino como predictores operativos del riesgo de anaerobiosis en el protocolo del Capítulo 5.

% ---------------------------------------------------------

\subsection{Nodo IoT y computación en el borde}
\label{subsec:cap2_nodo_iot}

El microcontrolador central del nodo es un ESP32-WROOM-32 (Espressif Systems), seleccionado por su bajo consumo energético (\SI{0.5}{\watt} en operación activa), conectividad dual Wi-Fi/Bluetooth y capacidad de procesamiento en el borde (dual-core \SI{240}{\mega\hertz}). Este dispositivo ejecuta tres tareas concurrentes: (i) adquisición multiparamétrica cada \SI{10}{\second}, (ii) filtrado digital de ruido mediante media móvil exponencial (EMA, $\alpha = 0.3$), y (iii) telemetría asíncrona mediante protocolo MQTT hacia un broker local (Raspberry Pi 4) o, en caso de disponibilidad, hacia un servidor remoto \parencite{tsaabitahDataCommunicationUsing2022, duobieneDevelopmentWirelessSensor2022}.

La arquitectura de comunicación sigue una topología de dos capas. La capa local (edge) almacena los datos en memoria Flash interna (\SI{4}{\mega\byte}) y en tarjeta microSD como respaldo físico, garantizando continuidad del registro ante fallos de red. La capa de agregación (gateway) ejecuta un broker MQTT ligero (Mosquitto) que recibe las tramas JSON codificadas con sello de tiempo (timestamp), identificador de réplica y checksum CRC16. Este diseño descentralizado permite operación autónoma en entornos rurales con conectividad intermitente, alineándose con los principios de computación en el borde y reducción de latencia para alertas operativas \parencite{yavariArtEMonArtificialIntelligence2023, rajRealTimeEstimation2023}.

El firmware del nodo se desarrolló en Arduino Framework con las siguientes rutinas de robustez: (a) watchdog timer de hardware que reinicia el sistema ante bloqueos mayores a \SI{30}{\second}; (b) validación de rangos físicos (sanitización) que descarta lecturas de \ce{CO2} $>$ \num{10000}~ppm o HR $>$ \SI{100}{\percent} como errores de sensor; y (c) sincronización horaria por NTP al inicio de cada ciclo de \SI{24}{\hour} para mantener trazabilidad temporal. La elección de componentes modulares (sensor desmontable, carcasa IP65 impresa en 3D) permite el reemplazo individual del elemento sensible sin descartar la unidad de procesamiento, reduciendo la generación de residuos electrónicos (e-waste) en un entorno corrosivo \parencite{vanIntegrationInternetofThingsSustainable2022}.
